\documentclass[11pt, a4paper]{article}
\usepackage[a4paper, margin=2.5cm]{geometry}
\usepackage{fontspec}
\usepackage[indonesian, bidi=basic, provide=*]{babel}
\babelprovide[import, onchar=ids fonts]{english}
\babelfont{rm}{Noto Sans}
\babelfont{tt}{Noto Sans Mono}
\usepackage{enumitem}
\setlist[itemize]{label=-}
\usepackage{amsmath, amssymb}
\usepackage{xcolor}
\usepackage{hyperref}

\hypersetup{
    colorlinks=true,
    linkcolor=blue,
    urlcolor=blue
}

\definecolor{emtcomment}{RGB}{0, 128, 0}
\definecolor{emtcommand}{RGB}{180, 0, 0}

\newcommand{\emtcmd}[1]{\noindent\texttt{\textcolor{emtcommand}{>}\ #1}\par\smallskip}

\begin{document}

\section*{Menggambar Grafik 2D dengan EMT}

Notebook ini menjelaskan tentang cara menggambar berbagai kurva dan grafik 2D dengan software EMT. EMT menyediakan fungsi \texttt{plot2d()} untuk menggambar berbagai kurva dan grafik dua dimensi (2D).

\subsection*{Plot Dasar}

Terdapat fungsi-fungsi dasar untuk membuat plot. Ada koordinat layar yang selalu berkisar dari 0 hingga 1024 pada setiap sumbunya, tidak peduli apakah layar berbentuk persegi atau tidak. Selain itu, ada koordinat plot yang dapat diatur dengan \texttt{setplot()}. Pemetaan di antara kedua koordinat tersebut bergantung pada jendela plot saat ini. Sebagai contoh, fungsi bawaan \texttt{shrinkwindow()} menyisakan ruang untuk label sumbu dan judul plot.

Pada contoh berikut, kita hanya menggambar beberapa garis acak dalam berbagai warna. Untuk detail mengenai fungsi-fungsi ini, pelajari fungsi inti EMT.

\emtcmd{clg; \textcolor{emtcomment}{// bersihkan layar}}
\emtcmd{window(0,0,1024,1024); \textcolor{emtcomment}{// gunakan seluruh area jendela}}
\emtcmd{setplot(0,1,0,1); \textcolor{emtcomment}{// atur koordinat plot}}
\emtcmd{hold on; \textcolor{emtcomment}{// mulai mode penimpaan (overwrite)}}
\emtcmd{n=100; X=random(n,2); Y=random(n,2); \textcolor{emtcomment}{// buat titik-titik acak}}
\emtcmd{colors=rgb(random(n),random(n),random(n)); \textcolor{emtcomment}{// buat warna-warna acak}}
\emtcmd{loop 1 to n; color(colors[\#]); plot(X[\#],Y[\#]); end; \textcolor{emtcomment}{// gambar plot}}
\emtcmd{hold off; \textcolor{emtcomment}{// akhiri mode penimpaan}}
\emtcmd{insimg; \textcolor{emtcomment}{// sisipkan ke dalam notebook}}
\emtcmd{reset;}

Perintah menahan grafik (\texttt{hold on}) diperlukan karena perintah \texttt{plot()} akan menghapus jendela plot sebelumnya. Untuk menghapus semua pengaturan yang telah kita buat, gunakan \texttt{reset()}.

Untuk menampilkan gambar hasil plot di layar notebook, perintah \texttt{plot2d()} dapat diakhiri dengan titik dua (\texttt{:}). Cara lain adalah perintah \texttt{plot2d()} diakhiri dengan titik koma (\texttt{;}), kemudian menggunakan perintah \texttt{insimg()} untuk menampilkan gambar hasil plot.

Sebagai contoh lain, kita menggambar sebuah plot sisipan (inset) di dalam plot lainnya. Hal ini dilakukan dengan mendefinisikan jendela plot yang lebih kecil. Perhatikan bahwa jendela ini tidak menyediakan ruang untuk label sumbu di luar jendela plot, sehingga kita perlu menambahkan margin seperlunya. Perhatikan pula bahwa kita menyimpan dan memulihkan jendela penuh, serta menahan plot saat ini ketika menggambar plot sisipan.

\emtcmd{plot2d("x\^{}3-x");}
\emtcmd{xw=200; yw=100; ww=300; hw=300;}
\emtcmd{ow=window();}
\emtcmd{window(xw,yw,xw+ww,yw+hw);}
\emtcmd{hold on;}
\emtcmd{barclear(xw-50,yw-10,ww+60,ww+60);}
\emtcmd{plot2d("x\^{}4-x",grid=6):}
\emtcmd{hold off;}
\emtcmd{window(ow);}

Plot dengan beberapa gambar (multi-figure) dapat dibuat dengan cara yang sama. Tersedia fungsi utilitas \texttt{figure()} untuk keperluan ini.

\subsection*{Aspek Plot}

Plot bawaan menggunakan jendela plot berbentuk persegi. Anda dapat mengubahnya dengan fungsi \texttt{aspect()}. Jangan lupa untuk mengatur ulang aspeknya setelah selesai. Anda juga dapat mengubah pengaturan bawaan ini melalui menu dengan memilih "Set Aspect" ke rasio aspek tertentu atau sesuai ukuran jendela grafik saat ini.

Namun, Anda juga dapat mengubahnya hanya untuk satu plot saja. Untuk itu, ukuran area plot saat ini diubah, dan jendela diatur agar label memiliki ruang yang cukup.

\emtcmd{aspect(2); \textcolor{emtcomment}{// rasio panjang dan lebar 2:1}}
\emtcmd{plot2d(["sin(x)","cos(x)"],0,2pi):}
\emtcmd{aspect();}
\emtcmd{reset;}

Fungsi \texttt{reset()} memulihkan pengaturan bawaan plot termasuk rasio aspeknya.

\section*{Plot 2D di Euler}

Euler Math Toolbox memiliki fitur plot 2D, baik untuk data maupun fungsi. EMT menggunakan fungsi \texttt{plot2d}. Fungsi ini dapat memplot fungsi dan data. Selain itu, dimungkinkan juga untuk membuat plot di Maxima menggunakan Gnuplot atau di Python menggunakan Matplotlib.

Euler dapat menggambar plot 2D dari:
\begin{itemize}
    \item ekspresi,
    \item fungsi, variabel, atau kurva berparameter,
    \item vektor nilai $x$-$y$,
    \item sebaran titik (clouds of points) pada bidang,
    \item kurva implisit dengan garis ketinggian (levels) atau daerah level,
    \item fungsi kompleks.
\end{itemize}

Gaya plot mencakup berbagai gaya untuk garis dan titik, diagram batang (bar plots), serta plot berarsir (shaded plots).

\section*{Plot Ekspresi atau Variabel}

Sebuah ekspresi tunggal dalam \texttt{"x"} (misalnya \texttt{"4*x\^{}2"}) atau nama sebuah fungsi (misalnya \texttt{"f"}) menghasilkan grafik dari fungsi tersebut. Berikut adalah contoh paling dasar yang menggunakan rentang bawaan dan mengatur rentang $y$ yang sesuai agar grafik termuat dengan baik.

Catatan: Jika Anda mengakhiri baris perintah dengan tanda titik dua \texttt{":"}, plot akan disisipkan ke dalam jendela teks. Jika tidak, tekan tombol \texttt{TAB} untuk melihat plot apabila jendela plot tertutup.

\emtcmd{plot2d("x\^{}2"):}
\emtcmd{aspect(1.5); plot2d("x\^{}3-x"):}
\emtcmd{a:=5.6; plot2d("exp(-a*x\^{}2)/a"); insimg(30); \textcolor{emtcomment}{// menampilkan gambar hasil plot setinggi 30 baris}}

Dari beberapa contoh sebelumnya Anda dapat melihat bahwa aslinya gambar plot menggunakan sumbu $X$ dengan rentang nilai dari $-2$ sampai dengan $2$. Untuk mengubah rentang nilai $X$ dan $Y$, Anda dapat menambahkan nilai-nilai batas $X$ (dan $Y$) di belakang ekspresi yang digambar.

Rentang plot diatur menggunakan parameter-parameter berikut:
\begin{itemize}
    \item \texttt{a,b}: rentang $x$ (bawaan $-2,2$)
    \item \texttt{c,d}: rentang $y$ (bawaan: menyesuaikan skala nilai)
    \item \texttt{r}: alternatif berupa jari-jari di sekitar pusat plot
    \item \texttt{cx,cy}: koordinat pusat plot (bawaan $0,0$)
\end{itemize}

\emtcmd{plot2d("x\^{}3-x",-1,2):}
\emtcmd{plot2d("sin(x)",-2*pi,2*pi): \textcolor{emtcomment}{// plot sin(x) pada interval [-2pi, 2pi]}}
\emtcmd{plot2d("cos(x)","sin(3*x)",xmin=0,xmax=2pi):}

Alternatif pengganti tanda titik dua adalah perintah \texttt{insimg(lines)}, yang menyisipkan plot dengan tinggi sejumlah baris teks yang ditentukan.

Pada menu opsi, plot dapat diatur agar muncul:
\begin{itemize}
    \item di jendela terpisah yang ukurannya dapat diubah,
    \item di dalam jendela notebook.
\end{itemize}

Gaya tampilan lainnya dapat diperoleh dengan perintah plot tertentu. Dalam kondisi apa pun, tekan tombol tabulator (\texttt{TAB}) untuk melihat plot jika tersembunyi.

Untuk membagi jendela menjadi beberapa plot, gunakan perintah \texttt{figure()}. Pada contoh berikut, kita memplot $x^1$ hingga $x^4$ ke dalam 4 bagian jendela. Perintah \texttt{figure(0)} mengatur ulang ke jendela bawaan.

\emtcmd{reset;}
\emtcmd{figure(2,2); ... \\
\hspace*{1em} for n=1 to 4; figure(n); plot2d("x\^{}"+n); end; ... \\
\hspace*{1em} figure(0):}

Pada \texttt{plot2d()}, tersedia gaya alternatif menggunakan parameter \texttt{grid=x}. Sebagai gambaran umum, kita tampilkan berbagai gaya kisi (grid) dalam satu gambar (lihat bagian bawah untuk penjelasan perintah \texttt{figure()}). Gaya \texttt{grid=0} tidak disertakan karena tidak menampilkan kisi maupun bingkai.

\emtcmd{figure(3,3); ... \\
\hspace*{1em} for k=1:9; figure(k); plot2d("x\^{}3-x",-2,1,grid=k); end; ... \\
\hspace*{1em} figure(0):}

Jika argumen pada \texttt{plot2d()} berupa sebuah ekspresi yang diikuti oleh empat bilangan, keempat bilangan tersebut merupakan rentang $x$ dan $y$ untuk plot. Sebagai alternatif, \texttt{a}, \texttt{b}, \texttt{c}, \texttt{d} dapat ditentukan sebagai parameter bernama seperti \texttt{a=...} dan seterusnya.

Pada contoh berikut, kita mengubah gaya kisi, menambahkan label, dan menggunakan label vertikal untuk sumbu $y$.

\emtcmd{aspect(1.5); plot2d("sin(x)",0,2pi,-1.2,1.2,grid=3,xl="x",yl="sin(x)"):}
\emtcmd{plot2d("sin(x)+cos(2*x)",0,4pi):}

Gambar yang dihasilkan dengan menyisipkan plot ke dalam jendela teks disimpan pada direktori yang sama dengan file notebook, secara bawaan berada di subdirektori bernama \texttt{"images"}. Gambar-gambar tersebut juga digunakan saat ekspor ke HTML. Anda cukup menandai gambar mana pun dan menyalinnya ke papan klip (clipboard) dengan \texttt{Ctrl-C}. Tentu saja, Anda juga dapat mengekspor grafik saat ini menggunakan fungsi pada menu File.

Fungsi atau ekspresi dalam \texttt{plot2d} dievaluasi secara adaptif. Untuk kecepatan lebih tinggi, matikan plot adaptif dengan \texttt{<adaptive} dan tentukan jumlah subinterval dengan \texttt{n=...} Hal ini hanya diperlukan pada kasus-kasus tertentu saja.

\emtcmd{plot2d("sign(x)*exp(-x\^{}2)",-1,1,<adaptive,n=10000):}
\emtcmd{plot2d("x\^{}x",r=1.2,cx=1,cy=1):}

Perhatikan bahwa $x^x$ tidak terdefinisi untuk $x \le 0$. Fungsi \texttt{plot2d} menangkap galat ini dan mulai menggambar segera setelah fungsinya terdefinisi. Hal ini berlaku untuk semua fungsi yang menghasilkan \texttt{NAN} di luar daerah definisinya.

\emtcmd{plot2d("log(x)",-0.1,2):}

Parameter \texttt{square=true} (atau \texttt{>square}) memilih rentang $y$ secara otomatis sehingga menghasilkan jendela plot berbentuk persegi. Perhatikan bahwa secara bawaan, Euler menggunakan ruang persegi di dalam jendela plot.

\emtcmd{plot2d("x\^{}3-x",>square):}
\emtcmd{plot2d(''integrate("sin(x)*exp(-x\^{}2)",0,x)'',0,2): \textcolor{emtcomment}{// plot integral}}

Jika Anda memerlukan ruang lebih luas untuk label $y$, panggil \texttt{shrinkwindow()} dengan parameter \texttt{smaller}, atau atur nilai positif untuk \texttt{"smaller"} di dalam \texttt{plot2d()}.

\emtcmd{plot2d("gamma(x)",1,10,yl="y-values",smaller=6,<vertical):}

Ekspresi simbolik juga dapat digunakan, karena disimpan sebagai ekspresi string sederhana.

\emtcmd{x=linspace(0,2pi,1000); plot2d(sin(5x),cos(7x)):}
\emtcmd{a:=5.6; expr \&= exp(-a*x\^{}2)/a; \textcolor{emtcomment}{// mendefinisikan ekspresi}}
\emtcmd{plot2d(expr,-2,2): \textcolor{emtcomment}{// plot dari -2 sampai 2}}
\emtcmd{plot2d(expr,r=1,thickness=2): \textcolor{emtcomment}{// plot dalam persegi di sekitar (0,0)}}
\emtcmd{plot2d(\&diff(expr,x),>add,style="--",color=red): \textcolor{emtcomment}{// menambahkan plot lain}}
\emtcmd{plot2d(\&diff(expr,x,2),a=-2,b=2,c=-2,d=1): \textcolor{emtcomment}{// plot dalam persegi panjang}}
\emtcmd{plot2d(\&diff(expr,x),a=-2,b=2,>square): \textcolor{emtcomment}{// menjaga plot tetap persegi}}
\emtcmd{plot2d("x\^{}2",0,1,steps=1,color=red,n=10):}
\emtcmd{plot2d("x\^{}2",>add,steps=2,color=blue,n=10):}

\section*{Fungsi Satu Parameter}

Fungsi pembuat plot terpenting untuk grafik bidang datar adalah \texttt{plot2d()}. Fungsi ini diimplementasikan dalam bahasa Euler pada berkas \texttt{"plot.e"}, yang dimuat saat program dimulai.

Berikut adalah beberapa contoh penggunaan fungsi. Seperti biasa di EMT, pada fungsi yang bekerja untuk fungsi atau ekspresi lain, Anda dapat meneruskan parameter tambahan (selain $x$) yang bukan merupakan variabel global ke fungsi tersebut menggunakan parameter titik koma atau koleksi pemanggilan (call collection).

\emtcmd{function f(x,a) := x\^{}2/a+a*x\^{}2-x; \textcolor{emtcomment}{// mendefinisikan fungsi}}
\emtcmd{a=0.3; plot2d("f",0,1;a): \textcolor{emtcomment}{// plot dengan a=0.3}}
\emtcmd{plot2d("f",0,1;0.4): \textcolor{emtcomment}{// plot dengan a=0.4}}
\emtcmd{plot2d(\{\{"f",0.2\}\},0,1): \textcolor{emtcomment}{// plot dengan a=0.2}}
\emtcmd{plot2d(\{\{"f(x,b)",b=0.1\}\},0,1): \textcolor{emtcomment}{// plot dengan b=0.1}}
\emtcmd{function f(x) := x\^{}3-x; ... \\
\hspace*{1em} plot2d("f",r=1):}

Berikut adalah ringkasan bentuk fungsi yang diterima:
\begin{itemize}
    \item ekspresi atau ekspresi simbolik dalam $x$,
    \item fungsi atau fungsi simbolik berdasarkan namanya sebagai string \texttt{"f"},
    \item fungsi simbolik cukup dengan menuliskan namanya \texttt{f}.
\end{itemize}

Fungsi \texttt{plot2d()} juga menerima fungsi simbolik. Untuk fungsi simbolik, cukup menuliskan namanya saja.

\emtcmd{function f(x) \&= diff(x\^{}x,x)}
$$x^x (\log(x) + 1)$$
\emtcmd{plot2d(f,0,2):}

Tentu saja, untuk ekspresi atau ekspresi simbolik, menuliskan nama variabelnya saja sudah cukup untuk memplotnya.

\emtcmd{expr \&= sin(x)*exp(-x)}
$$e^{-x} \sin(x)$$
\emtcmd{plot2d(expr,0,3pi):}
\emtcmd{function f(x) \&= x\^{}x;}
\emtcmd{plot2d(f,r=1,cx=1,cy=1,color=blue,thickness=2);}
\emtcmd{plot2d(\&diff(f(x),x),>add,color=red,style="-.-"):}

Untuk gaya garis tersedia berbagai pilihan:
\begin{itemize}
    \item \texttt{style="..."}. Pilih dari \texttt{"-"}, \texttt{"--"}, \texttt{"-."}, \texttt{"."}, \texttt{".-."}, \texttt{"-.-"}.
    \item \texttt{color}: Lihat di bawah untuk pilihan warna.
    \item \texttt{thickness}: Nilai bawaan adalah 1.
\end{itemize}

Warna dapat dipilih sebagai salah satu warna bawaan atau sebagai warna RGB:
\begin{itemize}
    \item \texttt{0..15}: indeks warna bawaan.
    \item konstanta warna: \texttt{white}, \texttt{black}, \texttt{red}, \texttt{green}, \texttt{blue}, \texttt{cyan}, \texttt{olive}, \texttt{lightgray}, \texttt{gray}, \texttt{darkgray}, \texttt{orange}, \texttt{lightgreen}, \texttt{turquoise}, \texttt{lightblue}, \texttt{lightorange}, \texttt{yellow}.
    \item \texttt{rgb(red,green,blue)}: parameter berupa bilangan riil dalam interval $[0,1]$.
\end{itemize}

\emtcmd{plot2d("exp(-x\^{}2)",r=2,color=red,thickness=3,style="--"):}

Berikut adalah tampilan warna-warna bawaan yang tersedia di EMT.
\emtcmd{aspect(2); columnsplot(ones(1,16),lab=0:15,grid=0,color=0:15):}

Namun, Anda juga dapat menggunakan warna apa saja.
\emtcmd{columnsplot(ones(1,16),grid=0,color=rgb(0,0,linspace(0,1,15))):}

\section*{Menggambar Beberapa Kurva pada Bidang Koordinat yang Sama}

Memplot lebih dari satu fungsi (fungsi majemuk) ke dalam satu jendela dapat dilakukan dengan beberapa cara. Salah satu caranya adalah menggunakan \texttt{>add} pada pemanggilan \texttt{plot2d} kedua dan seterusnya. Kita telah menggunakan fitur ini pada contoh-contoh di atas.

\emtcmd{aspect(); plot2d("cos(x)",r=2,grid=6); plot2d("x",style=".",>add):}
\emtcmd{aspect(1.5); plot2d("sin(x)",0,2pi); plot2d("cos(x)",color=blue,style="--",>add):}

Salah satu kegunaan \texttt{>add} adalah untuk menambahkan titik pada kurva.
\emtcmd{plot2d("sin(x)",0,pi); plot2d(2,sin(2),>points,>add):}

Kita tambahkan titik potong beserta labelnya (pada posisi \texttt{"cl"} yang berarti tengah kiri / center left), lalu menyisipkan hasilnya ke dalam notebook. Kita juga menambahkan judul pada plot.

\emtcmd{plot2d(["cos(x)","x"],r=1.1,cx=0.5,cy=0.5, ... \\
\hspace*{1em} color=[black,blue],style=["-","."], ... \\
\hspace*{1em} grid=1);}
\emtcmd{x0=solve("cos(x)-x",1); ... \\
\hspace*{1em} plot2d(x0,x0,>points,>add,title="Intersection Demo"); ... \\
\hspace*{1em} label("cos(x) = x",x0,x0,pos="cl",offset=20):}

Pada demonstrasi berikut, kita memplot fungsi $\mathrm{sinc}(x)=\sin(x)/x$ beserta ekspansi Taylor ke-8 dan ke-16. Kita menghitung ekspansi ini menggunakan Maxima melalui ekspresi simbolik. Plot ini dibuat dalam perintah multibaris berikut dengan tiga kali pemanggilan \texttt{plot2d()}. Pemanggilan kedua dan ketiga menggunakan flag \texttt{>add}, yang membuat plot menggunakan rentang sumbu sebelumnya. Kita tambahkan kotak label untuk menjelaskan fungsi-fungsi tersebut.

\emtcmd{\$taylor(sin(x)/x,x,0,4)}
\emtcmd{plot2d("sinc(x)",0,4pi,color=green,thickness=2); ... \\
\hspace*{1em} plot2d(\&taylor(sin(x)/x,x,0,8),>add,color=blue,style="--"); ... \\
\hspace*{1em} plot2d(\&taylor(sin(x)/x,x,0,16),>add,color=red,style="-.-"); ... \\
\hspace*{1em} labelbox(["sinc","T8","T16"],styles=["-","--","-.-"], ... \\
\hspace*{2em} colors=[black,blue,red]):}

Pada contoh berikut, kita membangkitkan Polinomial Bernstein:
$$B_i(x) = \binom{n}{i} x^i (1-x)^{n-i}$$

\emtcmd{plot2d("(1-x)\^{}10",0,1); \textcolor{emtcomment}{// plot fungsi pertama}}
\emtcmd{for i=1 to 10; plot2d("bin(10,i)*x\^{}i*(1-x)\^{}(10-i)",>add); end;}
\emtcmd{insimg;}

Metode kedua adalah menggunakan pasangan matriks nilai $x$ dan matriks nilai $y$ yang berukuran sama. Kita membangkitkan matriks nilai dengan satu Polinomial Bernstein pada setiap barisnya. Untuk itu, kita cukup menggunakan vektor kolom $i$. Pelajari pengantar bahasa matriks untuk mengetahui detail lebih lanjut.

\emtcmd{x=linspace(0,1,500);}
\emtcmd{n=10; k=(0:n)'; \textcolor{emtcomment}{// n adalah vektor baris, k adalah vektor kolom}}
\emtcmd{y=bin(n,k)*x\^{}k*(1-x)\^{}(n-k); \textcolor{emtcomment}{// y menjadi sebuah matriks}}
\emtcmd{plot2d(x,y):}

Perhatikan bahwa parameter warna dapat berupa vektor. Setiap warna akan digunakan untuk masing-masing baris pada matriks.
\emtcmd{x=linspace(0,1,200); y=x\^{}(1:10)'; plot2d(x,y,color=1:10):}

Metode lainnya adalah menggunakan vektor ekspresi (string). Anda kemudian dapat menggunakan larik warna, larik gaya garis, dan larik ketebalan garis dengan panjang yang sama.
\emtcmd{plot2d(["sin(x)","cos(x)"],0,2pi,color=4:5):}
\emtcmd{plot2d(["sin(x)","cos(x)"],0,2pi): \textcolor{emtcomment}{// plot vektor ekspresi}}

Kita dapat memperoleh vektor seperti itu dari Maxima menggunakan \texttt{makelist()} dan \texttt{mxm2str()}.
\emtcmd{v \&= makelist(binomial(10,i)*x\^{}i*(1-x)\^{}(10-i),i,0,10) \textcolor{emtcomment}{// membuat daftar ekspresi}}
$$\begin{aligned}
\big[&(1-x)^{10},\; 10(1-x)^9 x,\; 45(1-x)^8 x^2,\; 120(1-x)^7 x^3,\; 210(1-x)^6 x^4,\; 252(1-x)^5 x^5,\\
&210(1-x)^4 x^6,\; 120(1-x)^3 x^7,\; 45(1-x)^2 x^8,\; 10(1-x)x^9,\; x^{10}\big]
\end{aligned}$$

\emtcmd{mxm2str(v) \textcolor{emtcomment}{// mengubah vektor simbolik menjadi vektor string}}
\begin{verbatim}
    (1-x)^10
    10*(1-x)^9*x
    45*(1-x)^8*x^2
    120*(1-x)^7*x^3
    210*(1-x)^6*x^4
    252*(1-x)^5*x^5
    210*(1-x)^4*x^6
    120*(1-x)^3*x^7
    45*(1-x)^2*x^8
    10*(1-x)*x^9
    x^10
\end{verbatim}

\emtcmd{plot2d(mxm2str(v),0,1): \textcolor{emtcomment}{// memplot fungsi-fungsi}}

Alternatif lain adalah menggunakan bahasa matriks Euler. Jika suatu ekspresi menghasilkan matriks fungsi dengan satu fungsi pada setiap baris, semua fungsi tersebut akan digambar ke dalam satu plot. Untuk itu, gunakan vektor parameter dalam bentuk vektor kolom. Jika larik warna ditambahkan, warna tersebut akan digunakan untuk masing-masing baris pada plot.

\emtcmd{n=(1:10)'; plot2d("x\^{}n",0,1,color=1:10):}

Ekspresi dan fungsi satu baris dapat membaca variabel global. Jika Anda tidak dapat menggunakan variabel global, Anda perlu menggunakan fungsi dengan parameter tambahan dan meneruskan parameter tersebut sebagai parameter titik koma. Pastikan meletakkan semua parameter bernama di bagian akhir perintah \texttt{plot2d}. Pada contoh berikut, kita meneruskan $a=5$ ke fungsi $f$, yang kita plot dari $-10$ hingga $10$.

\emtcmd{function f(x,a) := 1/a*exp(-x\^{}2/a); ... \\
\hspace*{1em} plot2d("f",-10,10;5,thickness=2,title="a=5"):}

Sebagai alternatif, gunakan koleksi yang berisi nama fungsi beserta semua parameter tambahannya. Daftar khusus ini disebut koleksi pemanggilan (call collections), dan merupakan cara yang disarankan untuk meneruskan argumen ke suatu fungsi yang juga diteruskan sebagai argumen ke fungsi lain. Pada contoh berikut, kita menggunakan perulangan (loop) untuk memplot beberapa fungsi.

\emtcmd{plot2d(\{\{"f",1\}\},-10,10); ... \\
\hspace*{1em} for a=2:10; plot2d(\{\{"f",a\}\},>add); end:}

Kita dapat memperoleh hasil yang sama dengan cara berikut menggunakan bahasa matriks EMT. Setiap baris matriks $f(x,a)$ adalah satu fungsi. Selain itu, kita dapat mengatur warna untuk setiap baris matriks. Klik dua kali pada fungsi \texttt{getspectral()} untuk melihat penjelasannya.

\emtcmd{x=-10:0.01:10; a=(1:10)'; plot2d(x,f(x,a),color=getspectral(a/10)):}

\subsection*{Label Teks}

Dekorasi sederhana dapat berupa:
\begin{itemize}
    \item judul dengan \texttt{title="..."}
    \item label sumbu $x$ dan $y$ dengan \texttt{xl="..."}, \texttt{yl="..."}
    \item label teks lainnya dengan \texttt{label("...",x,y)}
\end{itemize}

Perintah \texttt{label} akan menggambar teks ke dalam plot saat ini pada koordinat plot $(x,y)$ dan dapat menerima argumen posisi.

\emtcmd{plot2d("x\^{}3-x",-1,2,title="y=x\^{}3-x",yl="y",xl="x"):}
\emtcmd{expr := "log(x)/x"; ... \\
\hspace*{1em} plot2d(expr,0.5,5,title="y="+expr,xl="x",yl="y"); ... \\
\hspace*{1em} label("(1,0)",1,0); label("Max",E,expr(E),pos="lc"):}

Tersedia pula fungsi \texttt{labelbox()} yang dapat menampilkan keterangan fungsi dan teks. Fungsi ini menerima vektor string dan warna, masing-masing satu item untuk setiap fungsi.

\emtcmd{function f(x) \&= x\^{}2*exp(-x\^{}2); ... \\
\hspace*{1em} plot2d(\&f(x),a=-3,b=3,c=-1,d=1); ... \\
\hspace*{1em} plot2d(\&diff(f(x),x),>add,color=blue,style="--"); ... \\
\hspace*{1em} labelbox(["function","derivative"],styles=["-","--"], ... \\
\hspace*{2em} colors=[black,blue],w=0.4):}

Secara bawaan kotak label ditempatkan di kanan atas, namun \texttt{>left} menempatkannya di kiri atas. Anda dapat memindahkannya ke posisi mana pun yang diinginkan. Posisi jangkarnya adalah sudut kanan atas kotak, dan angkanya merupakan pecahan dari ukuran jendela grafik. Lebarnya diatur secara otomatis.

Untuk plot titik, kotak label juga dapat digunakan. Tambahkan parameter \texttt{>points}, atau vektor flag, satu untuk setiap label. Pada contoh berikut hanya ada satu fungsi, sehingga kita dapat menggunakan string biasa alih-alih vektor string. Kita atur warna teks menjadi hitam untuk contoh ini.

\emtcmd{n=10; plot2d(0:n,bin(n,0:n),>addpoints); ... \\
\hspace*{1em} labelbox("Binomials",styles="[]",>points,x=0.1,y=0.1, ... \\
\hspace*{1em} tcolor=black,>left):}

Gaya plot seperti ini juga tersedia di \texttt{statplot()}. Seperti halnya \texttt{plot2d()}, warna dapat diatur untuk setiap baris plot. Tersedia berbagai plot khusus lainnya untuk keperluan statistika (lihat tutorial tentang statistika).

\emtcmd{statplot(1:10,random(2,10),color=[red,blue]):}

Fitur serupa adalah fungsi \texttt{textbox()}. Secara bawaan lebarnya mengikuti lebar maksimum baris teks, namun dapat pula diatur sendiri oleh pengguna.

\emtcmd{function f(x) \&= exp(-x)*sin(2*pi*x); ... \\
\hspace*{1em} plot2d("f(x)",0,2pi); ... \\
\hspace*{1em} textbox(latex("\textbackslash text\{Contoh osilasi teredam\}\textbackslash\ f(x)=e\^{}\{-x\}sin(2\textbackslash pi x)"),w=0.85):}

Label teks, judul, kotak label, dan teks lainnya dapat memuat string Unicode (lihat sintaks EMT untuk informasi lebih lanjut mengenai string Unicode).
\emtcmd{plot2d("x\^{}3-x",title=u"x \&rarr; x\&sup3; - x"):}

Label pada sumbu $x$ dan $y$ dapat dibuat vertikal, demikian pula sumbunya.
\emtcmd{plot2d("sinc(x)",0,2pi,xl="x",yl=u"x \&rarr; sinc(x)",>vertical):}

\subsection*{LaTeX}

Anda juga dapat memplot rumus LaTeX jika sistem LaTeX sudah terpasang di komputer Anda (disarankan MiKTeX). Jalur (path) ke berkas biner \texttt{"latex"} dan \texttt{"dvipng"} harus terdaftar pada system path, atau Anda perlu mengaturnya di menu opsi.

Perhatikan bahwa penguraian (parsing) LaTeX berjalan cukup lambat. Jika Anda ingin menggunakan LaTeX dalam plot animasi, sebaiknya panggil \texttt{latex()} satu kali sebelum perulangan dan gunakan hasilnya (berupa gambar dalam matriks RGB).

Pada plot berikut, kita menggunakan LaTeX untuk label sumbu $x$ dan $y$, sebuah label titik, kotak label, dan judul plot.
\emtcmd{plot2d("exp(-x)*sin(x)/x",a=0,b=2pi,c=0,d=1,grid=6,color=blue, ... \\
\hspace*{1em} title=latex("\textbackslash text\{Fungsi \$\textbackslash Phi\$\}"), ... \\
\hspace*{1em} xl=latex("\textbackslash phi"),yl=latex("\textbackslash Phi(\textbackslash phi)")); ... \\
\hspace*{1em} textbox( ... \\
\hspace*{2em} latex("\textbackslash Phi(\textbackslash phi) = e\^{}\{-\textbackslash phi\} \textbackslash frac\{\textbackslash sin(\textbackslash phi)\}\{\textbackslash phi\}"),x=0.8,y=0.5); ... \\
\hspace*{1em} label(latex("\textbackslash Phi",color=blue),1,0.4):}

Sering kali kita menginginkan jarak tanda skala yang tidak seragam beserta label teks pada sumbu $x$. Kita dapat menggunakan \texttt{xaxis()} dan \texttt{yaxis()} seperti yang akan ditunjukkan nanti.

Cara termudah adalah membuat plot kosong berbingkai menggunakan \texttt{grid=4}, lalu menambahkan garis kisi dengan \texttt{ygrid()} dan \texttt{xgrid()}. Pada contoh berikut, kita menggunakan tiga string LaTeX untuk label pada sumbu $x$ dengan \texttt{xtick()}.

\emtcmd{plot2d("sinc(x)",0,2pi,grid=4,<ticks); ... \\
\hspace*{1em} ygrid(-2:0.5:2,grid=6); ... \\
\hspace*{1em} xgrid([0:2]*pi,<ticks,grid=6); ... \\
\hspace*{1em} xtick([0,pi,2pi],["0","\textbackslash pi","2\textbackslash pi"],>latex):}

Tentu saja, fungsi juga dapat digunakan.
\emtcmd{function map f(x) ...}
\begin{verbatim}
    if x>0 then return x^4
    else return x^2
    endif
    endfunction
\end{verbatim}

Parameter \texttt{"map"} membantu agar fungsi dapat bekerja untuk vektor. Untuk sekadar menggambar plotnya hal ini sebenarnya tidak wajib, namun untuk menunjukkan manfaat vektorisasi, kita tambahkan beberapa titik penting pada plot di $x=-1$, $x=0$, dan $x=1$. Pada plot berikut, kita juga memasukkan beberapa kode LaTeX untuk dua label dan sebuah kotak teks. Tentu saja, fitur LaTeX hanya dapat digunakan jika LaTeX telah terpasang dengan benar.

\emtcmd{plot2d("f",-1,1,xl="x",yl="f(x)",grid=6); ... \\
\hspace*{1em} plot2d([-1,0,1],f([-1,0,1]),>points,>add); ... \\
\hspace*{1em} label(latex("x\^{}3"),0.72,f(0.72)); ... \\
\hspace*{1em} label(latex("x\^{}2"),-0.52,f(-0.52),pos="ll"); ... \\
\hspace*{1em} textbox( ... \\
\hspace*{2em} latex("f(x)=\textbackslash begin\{cases\} x\^{}3 \& x>0 \textbackslash\textbackslash\ x\^{}2 \& x \textbackslash le 0\textbackslash end\{cases\}"), ... \\
\hspace*{2em} x=0.7,y=0.2):}

\subsection*{Interaksi Pengguna}

Saat memplot fungsi atau ekspresi, parameter \texttt{>user} memungkinkan pengguna memperbesar (zoom) dan menggeser plot menggunakan tombol panah atau tetikus (mouse). Pengguna dapat:
\begin{itemize}
    \item memperbesar/memperkecil dengan tombol \texttt{+} atau \texttt{-}
    \item menggeser plot dengan tombol panah kursor
    \item memilih area jendela plot dengan tetikus
    \item mengatur ulang tampilan ke awal dengan tombol spasi
    \item keluar dengan tombol Enter (return)
\end{itemize}

Tombol spasi akan mengembalikan plot ke jendela plot semula. Saat memplot data, flag \texttt{>user} hanya akan menunggu penekanan tombol apa saja.

\emtcmd{plot2d(\{\{"x\^{}3-a*x",a=1\}\},>user,title="Tekan sembarang tombol!"):}
\emtcmd{plot2d("exp(x)*sin(x)",user=true, ... \\
\hspace*{1em} title="+/- atau tombol panah (Enter untuk keluar)"):}

Contoh berikut menunjukkan interaksi pengguna tingkat lanjut (lihat tutorial pemrograman untuk detail lebih lanjut). Fungsi bawaan \texttt{mousedrag()} menunggu aktivitas tetikus atau papan ketik, seperti penekanan tombol tetikus, pergeseran, maupun penekanan tombol ketik. Fungsi \texttt{dragpoints()} memanfaatkan fitur ini agar pengguna dapat menggeser titik mana pun pada plot.

Pertama kita memerlukan fungsi pembuat plot. Sebagai contoh, kita menginterpolasi 5 titik dengan polinomial. Fungsi ini harus memplot ke dalam area plot yang tetap.

\emtcmd{function plotf(xp,yp,select) ...}
\begin{verbatim}
      d=interp(xp,yp);
      plot2d("interpval(xp,d,x)";d,xp,r=2);
      plot2d(xp,yp,>points,>add);
      if select>0 then
        plot2d(xp[select],yp[select],color=red,>points,>add);
      endif;
      title("Geser satu titik, atau tekan spasi atau Enter!");
    endfunction
\end{verbatim}

Perhatikan parameter titik koma pada \texttt{plot2d} (\texttt{d} dan \texttt{xp}), yang diteruskan ke evaluasi fungsi \texttt{interp()}. Tanpa cara ini, kita harus menulis fungsi \texttt{plotinterp()} terlebih dahulu dengan mengakses nilai-nilai tersebut secara global.

Sekarang kita bangkitkan beberapa nilai acak, dan biarkan pengguna menggeser titik-titiknya.
\emtcmd{t=-1:0.5:1; dragpoints("plotf",t,random(size(t))-0.5):}

Tersedia juga fungsi yang memplot fungsi lain yang bergantung pada vektor parameter, dan memungkinkan pengguna menyesuaikan parameter tersebut secara interaktif. Pertama kita siapkan fungsi plotnya.
\emtcmd{function plotf([a,b]) := plot2d("exp(a*x)*cos(2pi*b*x)",0,2pi;a,b);}

Kemudian kita perlu menyiapkan nama parameter, nilai awal, matriks rentang berukuran $n \times 2$, serta baris judul opsional. Terdapat penggeser (slider) interaktif yang nilainya dapat diatur oleh pengguna melalui fungsi \texttt{dragvalues()}.

\emtcmd{dragvalues("plotf",["a","b"],[-1,2],[[-2,2];[1,10]], ... \\
\hspace*{1em} heading="Geser nilai-nilai berikut:",hcolor=black):}

Nilai yang digeser juga dapat dibatasi hanya pada bilangan bulat. Sebagai contoh, kita tulis fungsi yang memplot polinomial Taylor berderajat $n$ untuk fungsi kosinus.

\emtcmd{function plotf(n) ...}
\begin{verbatim}
    plot2d("cos(x)",0,2pi,>square,grid=6);
    plot2d(&"taylor(cos(x),x,0,@n)",color=blue,>add);
    textbox("Polinomial Taylor derajat "+n,0.1,0.02,style="t",>left);
    endfunction
\end{verbatim}

Sekarang kita biarkan derajat $n$ bervariasi dari 0 hingga 20 dalam 20 langkah. Hasil dari \texttt{dragvalues()} digunakan untuk menggambar sketsa dengan nilai $n$ tersebut, lalu menyisipkan plotnya ke dalam notebook.

\emtcmd{nd=dragvalues("plotf","degree",2,[0,20],20,y=0.8, ... \\
\hspace*{1em} heading="Geser nilai derajat:"); ... \\
\hspace*{1em} plotf(nd):}

Berikut adalah demonstrasi sederhana dari fungsi interaksi tersebut. Pengguna dapat menggambar langsung di atas jendela plot dan meninggalkan jejak titik-titik.

\emtcmd{function dragtest ...}
\begin{verbatim}
      plot2d(none,r=1,title="Geser dengan tetikus, atau tekan sembarang tombol!");
      start=0;
      repeat
        {flag,m,time}=mousedrag();
        if flag==0 then return; endif;
        if flag==2 then
          hold on; mark(m[1],m[2]); hold off;
        endif;
      end
    endfunction
\end{verbatim}

Tekan sembarang tombol untuk menghentikan animasi ini.
\emtcmd{dragtest \textcolor{emtcomment}{// lihat hasilnya dan cobalah lakukan!}}

\subsection*{Gaya Plot 2D}

Secara bawaan, EMT menghitung tanda skala sumbu secara otomatis dan menambahkan label pada setiap tanda skala. Hal ini dapat diubah dengan parameter \texttt{grid}. Gaya bawaan sumbu dan label dapat dimodifikasi, serta label dan judul dapat ditambahkan secara manual. Untuk kembali ke gaya bawaan, gunakan \texttt{reset()}.

\emtcmd{aspect();}
\emtcmd{figure(3,4); ... \\
\hspace*{1em} figure(1); plot2d("x\^{}3-x",grid=0); ... \textcolor{emtcomment}{// tanpa kisi, bingkai, maupun sumbu} \\
\hspace*{1em} figure(2); plot2d("x\^{}3-x",grid=1); ... \textcolor{emtcomment}{// hanya sumbu x dan y} \\
\hspace*{1em} figure(3); plot2d("x\^{}3-x",grid=2); ... \textcolor{emtcomment}{// tanda skala bawaan} \\
\hspace*{1em} figure(4); plot2d("x\^{}3-x",grid=3); ... \textcolor{emtcomment}{// sumbu x-y dengan label di dalam} \\
\hspace*{1em} figure(5); plot2d("x\^{}3-x",grid=4); ... \textcolor{emtcomment}{// tanpa tanda skala, hanya label} \\
\hspace*{1em} figure(6); plot2d("x\^{}3-x",grid=5); ... \textcolor{emtcomment}{// bawaan, namun tanpa margin} \\
\hspace*{1em} figure(7); plot2d("x\^{}3-x",grid=6); ... \textcolor{emtcomment}{// hanya sumbu} \\
\hspace*{1em} figure(8); plot2d("x\^{}3-x",grid=7); ... \textcolor{emtcomment}{// hanya sumbu, tanda skala di sumbu} \\
\hspace*{1em} figure(9); plot2d("x\^{}3-x",grid=8); ... \textcolor{emtcomment}{// hanya sumbu, sub-skala lebih halus di sumbu} \\
\hspace*{1em} figure(10); plot2d("x\^{}3-x",grid=9); ... \textcolor{emtcomment}{// bawaan, tanda skala kecil di dalam} \\
\hspace*{1em} figure(11); plot2d("x\^{}3-x",grid=10); ... \textcolor{emtcomment}{// tanpa tanda skala, hanya sumbu} \\
\hspace*{1em} figure(0):}

Parameter \texttt{<frame} mematikan bingkai, sedangkan \texttt{framecolor=blue} mengatur warna bingkai menjadi biru. Jika Anda ingin membuat tanda skala sendiri, gunakan \texttt{style=0} (atau \texttt{grid=0}), lalu tambahkan elemen lainnya secara manual.

\emtcmd{aspect(1.5);}
\emtcmd{plot2d("x\^{}3-x",grid=0); \textcolor{emtcomment}{// gambar plot}}
\emtcmd{frame; xgrid([-1,0,1]); ygrid(0): \textcolor{emtcomment}{// tambahkan bingkai dan garis kisi}}

Untuk judul plot dan label sumbu, perhatikan contoh berikut.
\emtcmd{plot2d("exp(x)",-1,1);}
\emtcmd{textcolor(black); \textcolor{emtcomment}{// atur warna teks menjadi hitam}}
\emtcmd{title(latex("y=e\^{}x")); \textcolor{emtcomment}{// judul di atas plot}}
\emtcmd{xlabel(latex("x")); \textcolor{emtcomment}{// "x" untuk sumbu x}}
\emtcmd{ylabel(latex("y"),>vertical); \textcolor{emtcomment}{// "y" vertikal untuk sumbu y}}
\emtcmd{label(latex("(0,1)"),0,1,color=blue): \textcolor{emtcomment}{// beri label pada titik}}

Sumbu koordinat dapat digambar secara terpisah dengan \texttt{xaxis()} dan \texttt{yaxis()}.
\emtcmd{plot2d("x\^{}3-x",<grid,<frame);}
\emtcmd{xaxis(0,xx=-2:1,style="->"); yaxis(0,yy=-5:5,style="->"):}

Teks pada plot dapat diatur dengan \texttt{label()}. Pada contoh berikut, \texttt{"lc"} berarti bawah tengah (lower center), yang mengatur posisi label relatif terhadap koordinat plot.

\emtcmd{function f(x) \&= x\^{}3-x}
$$x^3 - x$$
\emtcmd{plot2d(f,-1,1,>square);}
\emtcmd{x0=fmin(f,0,1); \textcolor{emtcomment}{// hitung titik minimum}}
\emtcmd{label("Rel. Min.",x0,f(x0),pos="lc"): \textcolor{emtcomment}{// tambahkan label di titik tersebut}}

Tersedia juga kotak teks (label box).
\emtcmd{plot2d(\&f(x),-1,1,-2,2); \textcolor{emtcomment}{// fungsi}}
\emtcmd{plot2d(\&diff(f(x),x),>add,style="--",color=red); \textcolor{emtcomment}{// turunan}}
\emtcmd{labelbox(["f","f'"],["-","--"],[black,red]): \textcolor{emtcomment}{// kotak label}}
\emtcmd{plot2d(["exp(x)","1+x"],color=[black,blue],style=["-","-.-"]):}
\emtcmd{gridstyle("->",color=gray,textcolor=gray,framecolor=gray); ... \\
\hspace*{1em} plot2d("x\^{}3-x",grid=1); ... \\
\hspace*{1em} settitle("y=x\^{}3-x",color=black); ... \\
\hspace*{1em} label("x",2,0,pos="bc",color=gray); ... \\
\hspace*{1em} label("y",0,6,pos="cl",color=gray); ... \\
\hspace*{1em} reset():}

Untuk kendali yang lebih leluasa, sumbu $x$ dan sumbu $y$ dapat dibuat secara manual. Perintah \texttt{fullwindow()} memperluas jendela plot karena kita tidak lagi membutuhkan ruang label di luar jendela plot. Gunakan \texttt{shrinkwindow()} atau \texttt{reset()} untuk kembali ke pengaturan bawaan.

\emtcmd{fullwindow; ... \\
\hspace*{1em} gridstyle(color=darkgray,textcolor=darkgray); ... \\
\hspace*{1em} plot2d(["2\^{}x","1","2\^{}(-x)"],a=-2,b=2,c=0,d=4,<grid,color=4:6,<frame); ... \\
\hspace*{1em} xaxis(0,-2:1,style="->"); xaxis(0,2,"x",<axis); ... \\
\hspace*{1em} yaxis(0,4,"y",style="->"); ... \\
\hspace*{1em} yaxis(-2,1:4,>left); ... \\
\hspace*{1em} yaxis(2,2\^{}(-2:2),style=".",<left); ... \\
\hspace*{1em} labelbox(["2\^{}x","1","2\^{}-x"],colors=4:6,x=0.8,y=0.2); ... \\
\hspace*{1em} reset:}

Berikut adalah contoh lain yang menggunakan string Unicode dan sumbu di luar area plot.
\emtcmd{aspect(1.5);}
\emtcmd{plot2d(["sin(x)","cos(x)"],0,2pi,color=[red,green],<grid,<frame); ... \\
\hspace*{1em} xaxis(-1.1,(0:2)*pi,xt=["0",u"\&pi;",u"2\&pi;"],style="-",>ticks,>zero); ... \\
\hspace*{1em} xgrid((0:0.5:2)*pi,<ticks); ... \\
\hspace*{1em} yaxis(-0.1*pi,-1:0.2:1,style="-",>zero,>grid); ... \\
\hspace*{1em} labelbox(["sin","cos"],colors=[red,green],x=0.5,y=0.2,>left); ... \\
\hspace*{1em} xlabel(u"\&phi;"); ylabel(u"f(\&phi;)"):}

\section*{Memplot Data 2D}

Jika $x$ dan $y$ adalah vektor data, data tersebut akan digunakan sebagai koordinat $x$ dan $y$ dari suatu kurva. Dalam hal ini, batas $a$, $b$, $c$, dan $d$, atau jari-jari $r$ dapat ditentukan, atau jendela plot akan menyesuaikan secara otomatis dengan data. Sebagai alternatif, \texttt{>square} dapat diatur untuk menjaga rasio aspek persegi.

Memplot sebuah ekspresi sebenarnya hanyalah bentuk singkat dari plot data. Untuk plot data, Anda memerlukan satu atau lebih baris nilai $x$, serta satu atau lebih baris nilai $y$. Dari rentang dan nilai $x$ tersebut, fungsi \texttt{plot2d} akan menghitung data yang akan diplot, secara bawaan dengan evaluasi adaptif. Untuk plot titik gunakan \texttt{">points"}, sedangkan untuk gabungan garis dan titik gunakan \texttt{">addpoints"}.

Namun, Anda juga dapat memasukkan data secara langsung:
\begin{itemize}
    \item Gunakan vektor baris untuk $x$ dan $y$ untuk menggambar satu fungsi.
    \item Matriks untuk $x$ dan $y$ akan diplot baris demi baris.
\end{itemize}

Berikut adalah contoh dengan satu baris untuk $x$ dan $y$.
\emtcmd{x=-10:0.1:10; y=exp(-x\^{}2)*x; plot2d(x,y):}

Data juga dapat diplot sebagai titik-titik. Gunakan \texttt{points=true} untuk keperluan ini. Plot bekerja seperti poligon, tetapi hanya menggambar titik-titik sudutnya saja:
\begin{itemize}
    \item \texttt{style="..."}: Pilih dari \texttt{"[]"}, \texttt{"<>"}, \texttt{"o"}, \texttt{"."}, \texttt{".."}, \texttt{"+"}, \texttt{"*"}, \texttt{"[]\#"}, \texttt{"<>\#"}, \texttt{"o\#"}, \texttt{"..\#"}, \texttt{"\#"}, \texttt{"|"}.
\end{itemize}

Untuk memplot himpunan titik, gunakan \texttt{>points}. Jika parameter warna berupa vektor warna, setiap titik akan mendapatkan warna yang berbeda. Untuk matriks koordinat dan vektor kolom warna, warna tersebut berlaku untuk setiap baris matriks. Parameter \texttt{>addpoints} menambahkan titik pada segmen garis untuk plot data.

\emtcmd{xdata=[1,1.5,2.5,3,4]; ydata=[3,3.1,2.8,2.9,2.7]; \textcolor{emtcomment}{// data}}
\emtcmd{plot2d(xdata,ydata,a=0.5,b=4.5,c=2.5,d=3.5,style="."); \textcolor{emtcomment}{// garis}}
\emtcmd{plot2d(xdata,ydata,>points,>add,style="o"): \textcolor{emtcomment}{// tambahkan titik}}
\emtcmd{p=polyfit(xdata,ydata,1); \textcolor{emtcomment}{// hitung garis regresi}}
\emtcmd{plot2d("polyval(p,x)",>add,color=red): \textcolor{emtcomment}{// tambahkan plot garis regresi}}

\section*{Menggambar Daerah yang Dibatasi Kurva}

Plot data pada dasarnya adalah poligon. Kita juga dapat memplot kurva atau kurva berisi (terarsir):
\begin{itemize}
    \item \texttt{filled=true} mengisi bagian dalam plot.
    \item \texttt{style="..."}: Pilih dari \texttt{"\#"}, \texttt{"/"}, \texttt{"\textbackslash"}, \texttt{"\textbackslash/"}.
    \item \texttt{fillcolor}: Lihat bagian sebelumnya untuk daftar warna yang tersedia.
\end{itemize}

Warna isian ditentukan oleh argumen \texttt{"fillcolor"}, dan parameter opsional \texttt{<outline} mencegah penggambaran garis batas untuk semua gaya kecuali gaya bawaan.

\emtcmd{t=linspace(0,2pi,1000); \textcolor{emtcomment}{// parameter untuk kurva}}
\emtcmd{x=sin(t)*exp(t/pi); y=cos(t)*exp(t/pi); \textcolor{emtcomment}{// x(t) dan y(t)}}
\emtcmd{figure(1,2); aspect(16/9)}
\emtcmd{figure(1); plot2d(x,y,r=10); \textcolor{emtcomment}{// plot kurva}}
\emtcmd{figure(2); plot2d(x,y,r=10,>filled,style="/",fillcolor=red); \textcolor{emtcomment}{// isi kurva}}
\emtcmd{figure(0):}

Pada contoh-contoh berikut kita memplot sebuah elips terisi dan dua segi enam terisi menggunakan kurva tertutup dengan 6 titik serta gaya arsiran yang berbeda.
\emtcmd{x=linspace(0,2pi,1000); plot2d(sin(x),cos(x)*0.5,r=1,>filled,style="/"):}
\emtcmd{t=linspace(0,2pi,6); ... \\
\hspace*{1em} plot2d(cos(t),sin(t),>filled,style="/",fillcolor=red,r=1.2):}
\emtcmd{t=linspace(0,2pi,6); plot2d(cos(t),sin(t),>filled,style="\#"):}

Contoh lainnya adalah segi tujuh (septagon), yang kita buat menggunakan 7 titik pada lingkaran satuan.
\emtcmd{t=linspace(0,2pi,7); ... \\
\hspace*{1em} plot2d(cos(t),sin(t),r=1,>filled,style="/",fillcolor=red):}

Berikut adalah himpunan nilai maksimum dari empat pertidaksamaan linear yang kurang dari atau sama dengan 3, yakni $A[k] \cdot v \le 3$ untuk semua baris matriks $A$. Agar sudut-sudutnya terlihat tajam, kita gunakan nilai $n$ yang cukup besar.
\emtcmd{A=[2,1;1,2;-1,0;0,-1];}
\emtcmd{function f(x,y) := max([x,y].A');}
\emtcmd{plot2d("f",r=4,level=[0;3],color=green,n=111):}

Keunggulan utama bahasa matriks adalah kemudahannya dalam membangkitkan tabel nilai fungsi.
\emtcmd{t=linspace(0,2pi,1000); x=cos(3*t); y=sin(4*t);}

Sekarang kita memiliki vektor nilai $x$ dan $y$. Fungsi \texttt{plot2d()} dapat memplot nilai-nilai ini sebagai kurva yang menghubungkan titik-titik tersebut, dan plotnya dapat diisi. Pada kasus ini dihasilkan pola arsiran yang menarik berkat aturan lilitan (winding rule) yang digunakan untuk pengisian area.
\emtcmd{plot2d(x,y,<grid,<frame,>filled):}

Sebuah vektor interval dapat diplot terhadap nilai $x$ sebagai daerah terisi di antara batas bawah dan batas atas interval tersebut. Fitur ini berguna untuk memplot galat komputasi maupun galat statistik.
\emtcmd{t=0:0.1:1; ... \\
\hspace*{1em} plot2d(t,interval(t-random(size(t)),t+random(size(t))),style="|"); ... \\
\hspace*{1em} plot2d(t,t,add=true):}

Jika $x$ adalah vektor terurut dan $y$ adalah vektor interval, maka \texttt{plot2d} akan memplot daerah terisi dari interval-interval tersebut pada bidang. Gaya pengisiannya sama dengan gaya pada poligon.
\emtcmd{t=-1:0.01:1; x=\~t-0.01,t+0.01\~; y=x\^{}3-x;}
\emtcmd{plot2d(t,y):}

Dimungkinkan juga untuk mengarsir daerah nilai untuk suatu fungsi tertentu. Untuk itu, \texttt{level} harus berupa matriks $2 \times n$, dengan baris pertama berisi batas bawah dan baris kedua berisi batas atas.
\emtcmd{expr := "2*x\^{}2+x*y+3*y\^{}4+y"; \textcolor{emtcomment}{// mendefinisikan ekspresi f(x,y)}}
\emtcmd{plot2d(expr,level=[0;1],style="-",color=blue): \textcolor{emtcomment}{// 0 <= f(x,y) <= 1}}

Kita juga dapat mengarsir rentang nilai seperti
$$-1 \le (x^2+y^2)^2-x^2+y^2 \le 0.$$
\emtcmd{plot2d("(x\^{}2+y\^{}2)\^{}2-x\^{}2+y\^{}2",r=1.2,level=[-1;0],style="/"):}
\emtcmd{plot2d("cos(x)","sin(x)\^{}3",xmin=0,xmax=2pi,>filled,style="/"):}

\section*{Grafik Fungsi Parametrik}

Nilai $x$ tidak harus terurut. Pasangan $(x,y)$ secara umum menggambarkan sebuah kurva; jika $x$ terurut, kurva tersebut merupakan grafik fungsi. Pada contoh berikut, kita memplot spiral
$$\gamma(t) = t \cdot (\cos(2\pi t),\sin(2\pi t))$$

Kita perlu menggunakan titik dalam jumlah banyak agar kurva terlihat mulus, atau menggunakan fungsi \texttt{adaptive()} untuk mengevaluasi ekspresi tersebut.
\emtcmd{t=linspace(0,1,1000); ... \\
\hspace*{1em} plot2d(t*cos(2*pi*t),t*sin(2*pi*t),r=1):}

Sebagai alternatif, dimungkinkan menggunakan dua ekspresi untuk menggambar kurva. Perintah berikut memplot kurva yang sama seperti di atas.
\emtcmd{plot2d("x*cos(2*pi*x)","x*sin(2*pi*x)",xmin=0,xmax=1,r=1):}
\emtcmd{t=linspace(0,1,1000); r=exp(-t); x=r*cos(2pi*t); y=r*sin(2pi*t);}
\emtcmd{plot2d(x,y,r=1):}

Pada contoh berikutnya, kita memplot kurva
$$\gamma(t) = (r(t) \cos(t), r(t) \sin(t))$$
dengan
$$r(t) = 1 + \frac{\sin(3t)}{2}.$$
\emtcmd{t=linspace(0,2pi,1000); r=1+sin(3*t)/2; x=r*cos(t); y=r*sin(t); ... \\
\hspace*{1em} plot2d(x,y,>filled,fillcolor=red,style="/",r=1.5):}

\section*{Menggambar Grafik Bilangan Kompleks}

Larik bilangan kompleks juga dapat diplot sehingga titik-titik kisinya saling terhubung. Jika jumlah garis kisi ditentukan (atau berupa vektor $1 \times 2$) pada argumen \texttt{cgrid}, hanya garis kisi tersebut yang akan ditampilkan. Matriks bilangan kompleks secara otomatis akan diplot sebagai kisi pada bidang kompleks.

Pada contoh berikut, kita memplot bayangan lingkaran satuan di bawah fungsi eksponensial. Parameter \texttt{cgrid} menyembunyikan sebagian kurva kisi.
\emtcmd{aspect(); r=linspace(0,1,50); a=linspace(0,2pi,80)'; z=r*exp(I*a); ... \\
\hspace*{1em} plot2d(z,a=-1.25,b=1.25,c=-1.25,d=1.25,cgrid=10):}
\emtcmd{aspect(1.25); r=linspace(0,1,50); a=linspace(0,2pi,200)'; z=r*exp(I*a);}
\emtcmd{plot2d(exp(z),cgrid=[40,10]):}
\emtcmd{r=linspace(0,1,10); a=linspace(0,2pi,40)'; z=r*exp(I*a);}
\emtcmd{plot2d(exp(z),>points,>add):}

Vektor bilangan kompleks secara otomatis diplot sebagai kurva pada bidang kompleks dengan sumbu bagian riil dan bagian imajiner. Pada contoh ini, kita memplot kurva dengan
$$\gamma(t) = e^{it} + e^{4it}$$
\emtcmd{t=linspace(0,2pi,1000); ... \\
\hspace*{1em} plot2d(exp(I*t)+exp(4*I*t),r=2):}

\section*{Plot Statistika}

Terdapat banyak fungsi yang dikhususkan untuk plot statistika. Salah satu plot yang sering digunakan adalah diagram kolom (column plot). Jumlah kumulatif dari nilai berdistribusi normal baku menghasilkan jalan acak (random walk).
\emtcmd{plot2d(cumsum(randnormal(1,1000))):}

Menggunakan dua baris menunjukkan jalan acak dalam dua dimensi.
\emtcmd{X=cumsum(randnormal(2,1000)); plot2d(X[1],X[2]):}
\emtcmd{columnsplot(cumsum(random(10)),style="/",color=blue):}

Diagram ini juga dapat menampilkan string sebagai label.
\emtcmd{months=["Jan","Feb","Mar","Apr","May","Jun", ... \\
\hspace*{1em} "Jul","Aug","Sep","Oct","Nov","Dec"];}
\emtcmd{values=[10,12,12,18,22,28,30,26,22,18,12,8];}
\emtcmd{columnsplot(values,lab=months,color=red,style="-");}
\emtcmd{title("Temperature"):}
\emtcmd{k=0:10;}
\emtcmd{plot2d(k,bin(10,k),>bar):}
\emtcmd{plot2d(k,bin(10,k)); plot2d(k,bin(10,k),>points,>add):}
\emtcmd{plot2d(normal(1000),normal(1000),>points,grid=6,style=".."):}
\emtcmd{plot2d(normal(1,1000),>distribution,style="O"):}
\emtcmd{plot2d("qnormal",0,5;2.5,0.5,>filled):}

Untuk memplot distribusi statistik hasil percobaan, Anda dapat menggunakan \texttt{distribution=n} pada \texttt{plot2d}.
\emtcmd{w=randexponential(1,1000); \textcolor{emtcomment}{// distribusi eksponensial}}
\emtcmd{plot2d(w,>distribution): \textcolor{emtcomment}{// atau distribution=n dengan n interval}}

Atau Anda dapat menghitung distribusinya dari data dan memplot hasilnya dengan \texttt{>bar} pada \texttt{plot3d}, atau menggunakan diagram kolom.
\emtcmd{w=normal(1000); \textcolor{emtcomment}{// distribusi normal baku 0-1}}
\emtcmd{\{x,y\}=histo(w,10,v=[-6,-4,-2,-1,0,1,2,4,6]); \textcolor{emtcomment}{// batas-batas interval v}}
\emtcmd{plot2d(x,y,>bar):}

Fungsi \texttt{statplot()} mengatur gaya plot menggunakan string sederhana.
\emtcmd{statplot(1:10,cumsum(random(10)),"b"):}
\emtcmd{n=10; i=0:n; ... \\
\hspace*{1em} plot2d(i,bin(n,i)/2\^{}n,a=0,b=10,c=0,d=0.3); ... \\
\hspace*{1em} plot2d(i,bin(n,i)/2\^{}n,points=true,style="ow",add=true,color=blue):}

Selain itu, data dapat diplot sebagai diagram batang (bar). Dalam hal ini, $x$ harus terurut dan panjangnya satu elemen lebih banyak daripada $y$. Batang akan membentang dari $x[i]$ hingga $x[i+1]$ dengan tinggi $y[i]$. Jika ukuran $x$ sama dengan $y$, vektor $x$ akan diperpanjang satu elemen mengikuti jarak terakhirnya. Gaya pengisian dapat digunakan seperti sebelumnya.

\emtcmd{n=10; k=bin(n,0:n); ... \\
\hspace*{1em} plot2d(-0.5:n+0.5,k,bar=true,fillcolor=lightgray):}

Data untuk diagram batang (\texttt{bar=1}) dan histogram (\texttt{histogram=1}) dapat diberikan secara eksplisit dalam \texttt{xv} dan \texttt{yv}, atau dihitung dari distribusi empiris pada \texttt{xv} dengan \texttt{>distribution} (atau \texttt{distribution=n}). Histogram dari nilai \texttt{xv} akan dihitung otomatis dengan \texttt{>histogram}. Jika \texttt{>even} ditentukan, nilai \texttt{xv} akan dihitung dalam interval bilangan bulat.

\emtcmd{plot2d(normal(10000),distribution=50):}
\emtcmd{k=0:10; m=bin(10,k); x=(0:11)-0.5; plot2d(x,m,>bar):}
\emtcmd{columnsplot(m,k):}
\emtcmd{plot2d(random(600)*6,histogram=6):}

Untuk distribusi, tersedia parameter \texttt{distribution=n}, yang menghitung nilai secara otomatis dan menampilkan distribusi relatif dengan $n$ subinterval.
\emtcmd{plot2d(normal(1,1000),distribution=10,style="\textbackslash/"):}

Dengan parameter \texttt{even=true}, perhitungan akan menggunakan interval bilangan bulat.
\emtcmd{plot2d(intrandom(1,1000,10),distribution=10,even=true):}

Perhatikan bahwa terdapat banyak plot statistika lain yang bermanfaat. Silakan lihat tutorial mengenai statistika.
\emtcmd{columnsplot(getmultiplicities(1:6,intrandom(1,6000,6))):}
\emtcmd{plot2d(normal(1,1000),>distribution); ... \\
\hspace*{1em} plot2d("qnormal(x)",color=red,thickness=2,>add):}

Tersedia pula plot khusus lainnya untuk statistika. Diagram kotak garis (boxplot) menampilkan kuartil distribusi beserta pencilan (outliers). Menurut definisi, pencilan dalam boxplot adalah data yang melampaui 1,5 kali jangkauan interkuartil (50\% bagian tengah data).
\emtcmd{M=normal(5,1000); boxplot(quartiles(M)):}

\section*{Fungsi Implisit}

Plot implisit menampilkan garis ketinggian (level lines) yang memenuhi $f(x,y)=\text{level}$, di mana \texttt{"level"} dapat berupa nilai tunggal maupun vektor nilai. Jika \texttt{level="auto"}, akan dibuat sebanyak \texttt{nc} garis level yang tersebar merata di antara nilai minimum dan maksimum fungsi. Gradasi warna lebih gelap atau lebih terang dapat ditambahkan dengan \texttt{>hue} untuk menunjukkan nilai fungsi. Untuk fungsi implisit, \texttt{xv} harus berupa fungsi atau ekspresi dari parameter $x$ dan $y$, atau dapat pula berupa matriks nilai.

Euler dapat menandai garis-garis level
$$f(x,y) = c$$
dari sembarang fungsi. Untuk menggambar himpunan $f(x,y)=c$ bagi satu atau lebih konstanta $c$, gunakan \texttt{plot2d()} dengan parameter \texttt{level=c}, di mana $c$ dapat berupa vektor garis level. Selain itu, skema warna dapat digambar pada latar belakang untuk menunjukkan nilai fungsi di setiap titik plot. Parameter \texttt{"n"} menentukan kehalusan plot.

\emtcmd{aspect(1.5);}
\emtcmd{plot2d("x\^{}2+y\^{}2-x*y-x",r=1.5,level=0,contourcolor=red):}
\emtcmd{expr := "2*x\^{}2+x*y+3*y\^{}4+y"; \textcolor{emtcomment}{// mendefinisikan ekspresi f(x,y)}}
\emtcmd{plot2d(expr,level=0): \textcolor{emtcomment}{// solusi dari f(x,y)=0}}
\emtcmd{plot2d(expr,level=0:0.5:20,>hue,contourcolor=white,n=200): \textcolor{emtcomment}{// tampilan bagus}}
\emtcmd{plot2d(expr,level=0:0.5:20,>hue,>spectral,n=200,grid=4): \textcolor{emtcomment}{// tampilan lebih bagus}}

Cara ini juga berlaku untuk plot data, namun Anda harus menentukan rentang untuk label sumbunya.
\emtcmd{x=-2:0.05:1; y=x'; z=expr(x,y);}
\emtcmd{plot2d(z,level=0,a=-1,b=2,c=-2,d=1,>hue):}
\emtcmd{plot2d("x\^{}3-y\^{}2",>contour,>hue,>spectral):}
\emtcmd{plot2d("x\^{}3-y\^{}2",level=0,contourwidth=3,>add,contourcolor=red):}
\emtcmd{z=z+normal(size(z))*0.2;}
\emtcmd{plot2d(z,level=0.5,a=-1,b=2,c=-2,d=1):}
\emtcmd{plot2d(expr,level=[0:0.2:5;0.05:0.2:5.05],color=lightgray):}
\emtcmd{plot2d("x\^{}2+y\^{}3+x*y",level=1,r=4,n=100):}
\emtcmd{plot2d("x\^{}2+2*y\^{}2-x*y",level=0:0.1:10,n=100,contourcolor=white,>hue):}

Dimungkinkan pula untuk mengarsir himpunan
$$a \le f(x,y) \le b$$
dengan rentang level. Untuk mengarsir daerah nilai bagi fungsi tertentu, \texttt{level} harus berupa matriks $2 \times n$ dengan baris pertama berisi batas bawah dan baris kedua berisi batas atas.

\emtcmd{plot2d(expr,level=[0;1],style="-",color=blue): \textcolor{emtcomment}{// 0 <= f(x,y) <= 1}}

Plot implisit juga dapat menampilkan rentang level. Parameter \texttt{level} harus berupa matriks interval level berukuran $2 \times n$, di mana baris pertama berisi nilai awal dan baris kedua berisi nilai akhir setiap interval. Sebagai alternatif, vektor baris sederhana dapat digunakan untuk \texttt{level}, dan parameter \texttt{dl} memperluas nilai level tersebut menjadi interval.

\emtcmd{plot2d("x\^{}4+y\^{}4",r=1.5,level=[0;1],color=blue,style="/"):}
\emtcmd{plot2d("x\^{}2+y\^{}3+x*y",level=[0,2,4;1,3,5],style="/",r=2,n=100):}
\emtcmd{plot2d("x\^{}2+y\^{}3+x*y",level=-10:20,r=2,style="-",dl=0.1,n=100):}
\emtcmd{plot2d("sin(x)*cos(y)",r=pi,>hue,>levels,n=100):}

Kita juga dapat menandai suatu daerah
$$a \le f(x,y) \le b$$
dengan menambahkan matriks level dua baris.
\emtcmd{plot2d("(x\^{}2+y\^{}2-1)\^{}3-x\^{}2*y\^{}3",r=1.3, ... \\
\hspace*{1em} style="\#",color=red,<outline, ... \\
\hspace*{1em} level=[-2;0],n=100):}

Kita dapat menentukan nilai level tertentu, misalnya untuk memplot solusi persamaan seperti
$$x^3-xy+x^2y^2=6$$
\emtcmd{plot2d("x\^{}3-x*y+x\^{}2*y\^{}2",r=6,level=1,n=100):}

\emtcmd{function starplot1 (v, style="/", color=green, lab=none) ...}
\begin{verbatim}
      if !holding() then clg; endif;
      w=window(); window(0,0,1024,1024);
      h=holding(1);
      r=max(abs(v))*1.2;
      setplot(-r,r,-r,r);
      n=cols(v); t=linspace(0,2pi,n);
      v=v|v[1]; c=v*cos(t); s=v*sin(t);
      cl=barcolor(color); st=barstyle(style);
      loop 1 to n
        polygon([0,c[#],c[#+1]],[0,s[#],s[#+1]],1);
        if lab!=none then
          rlab=v[#]+r*0.1;
          {col,row}=toscreen(cos(t[#])*rlab,sin(t[#])*rlab);
          ctext(""+lab[#],col,row-textheight()/2);
        endif;
      end;
      barcolor(cl); barstyle(st);
      holding(h);
      window(w);
    endfunction
\end{verbatim}

Di sini tidak digunakan garis kisi maupun tanda skala sumbu, dan kita memakai seluruh area jendela untuk plot. Kita memanggil \texttt{reset} sebelum menguji plot ini untuk memulihkan pengaturan grafik bawaan.

\emtcmd{reset; starplot1(normal(1,10)+5,color=red,lab=1:10):}

Terkadang Anda ingin menggambar sesuatu yang hampir dapat dilakukan oleh \texttt{plot2d}, tetapi tidak secara langsung. Pada fungsi berikut, kita membuat plot impuls logaritmik. Fungsi \texttt{plot2d} dapat membuat plot logaritmik, namun tidak untuk batang impuls.

\emtcmd{function logimpulseplot1 (x,y) ...}
\begin{verbatim}
      {x0,y0}=makeimpulse(x,log(y)/log(10));
      plot2d(x0,y0,>bar,grid=0);
      h=holding(1);
      frame();
      xgrid(ticks(x));
      p=plot();
      for i=-10 to 10;
        if i<=p[4] and i>=p[3] then
           ygrid(i,yt="10^"+i);
        endif;
      end;
      holding(h);
    endfunction
\end{verbatim}

Mari kita uji fungsi tersebut dengan nilai-nilai yang berdistribusi eksponensial.
\emtcmd{aspect(1.5); x=1:10; y=-log(random(size(x)))*200; ... \\
\hspace*{1em} logimpulseplot1(x,y):}

Mari kita buat animasi kurva 2D menggunakan plot langsung. Perintah \texttt{plot(x,y)} menggambar kurva ke dalam jendela plot, sedangkan \texttt{setplot(a,b,c,d)} mengatur jendela tersebut. Fungsi \texttt{wait(0)} memaksa plot segera tampil pada jendela grafik; tanpa perintah ini, penggambaran ulang hanya terjadi pada selang waktu tertentu.

\emtcmd{function animliss (n,m) ...}
\begin{verbatim}
    t=linspace(0,2pi,500);
    f=0;
    c=framecolor(0);
    l=linewidth(2);
    setplot(-1,1,-1,1);
    repeat
      clg;
      plot(sin(n*t),cos(m*t+f));
      wait(0);
      if testkey() then break; endif;
      f=f+0.02;
    end;
    framecolor(c);
    linewidth(l);
    endfunction
\end{verbatim}

Tekan sembarang tombol untuk menghentikan animasi ini.
\emtcmd{animliss(2,3); \textcolor{emtcomment}{// lihat hasilnya, jika sudah puas, tekan ENTER}}

\section*{Plot Logaritmik}

EMT menggunakan parameter \texttt{"logplot"} untuk skala logaritmik. Plot logaritmik dapat digambar menggunakan skala logaritmik pada sumbu $y$ dengan \texttt{logplot=1}, pada sumbu $x$ dan $y$ dengan \texttt{logplot=2}, atau pada sumbu $x$ saja dengan \texttt{logplot=3}.

\begin{itemize}
    \item \texttt{logplot=1}: logaritmik pada sumbu $y$
    \item \texttt{logplot=2}: logaritmik pada sumbu $x$ dan $y$
    \item \texttt{logplot=3}: logaritmik pada sumbu $x$
\end{itemize}

\emtcmd{plot2d("exp(x\^{}3-x)*x\^{}2",1,5,logplot=1):}
\emtcmd{plot2d("exp(x+sin(x))",0,100,logplot=1):}
\emtcmd{plot2d("exp(x+sin(x))",10,100,logplot=2):}
\emtcmd{plot2d("gamma(x)",1,10,logplot=1):}
\emtcmd{plot2d("log(x*(2+sin(x/100)))",10,1000,logplot=3):}

Fitur ini juga dapat digunakan untuk plot data.
\emtcmd{x=10\^{}(1:20); y=x\^{}2-x;}
\emtcmd{plot2d(x,y,logplot=2):}

\section*{Rujukan Lengkap Fungsi \texttt{plot2d()}}

\begin{verbatim}
function plot2d (xv, yv, btest, a, b, c, d, xmin, xmax, r, n, ..
  logplot, grid, frame, framecolor, square, color, thickness, style, ..
  auto, add, user, delta, points, addpoints, pointstyle, bar, histogram, ..
  distribution, even, steps, own, adaptive, hue, level, contour, ..
  nc, filled, fillcolor, outline, title, xl, yl, maps, contourcolor, ..
  contourwidth, ticks, margin, clipping, cx, cy, insimg, spectral, ..
  cgrid, vertical, smaller, dl, niveau, levels)
\end{verbatim}

Fungsi plot serbaguna untuk membuat grafik pada bidang datar (plot 2D). Fungsi ini dapat menggambar plot fungsi satu variabel, plot data, kurva pada bidang, diagram batang, kisi bilangan kompleks, serta plot implisit dari fungsi dua variabel.

\subsection*{Parameter}

\begin{verbatim}
x,y       : persamaan, fungsi, atau vektor data
a,b,c,d   : Area plot (bawaan a=-2, b=2)
r         : jika r diatur, maka a=cx-r, b=cx+r, c=cy-r, d=cy+r
            r dapat berupa vektor [rx,ry] atau [rx1,rx2,ry1,ry2].
xmin,xmax : rentang parameter untuk kurva
auto      : Menentukan rentang y secara otomatis (bawaan)
square    : jika true, menjaga rentang x-y tetap berbentuk persegi
n         : jumlah interval (bawaan bersifat adaptif)
grid      : 0 = tanpa kisi dan label,
            1 = hanya sumbu,
            2 = kisi normal (lihat penjelasan jumlah garis kisi)
            3 = sumbu di bagian dalam
            4 = tanpa kisi
            5 = kisi penuh termasuk margin
            6 = tanda skala pada bingkai
            7 = hanya sumbu
            8 = hanya sumbu dengan sub-skala
frame     : 0 = tanpa bingkai
framecolor: warna bingkai dan kisi
margin    : bilangan antara 0 dan 0.4 untuk margin di sekeliling plot
color     : Warna kurva. Jika berupa vektor warna, akan digunakan untuk
            setiap baris pada matriks plot. Untuk plot titik, sebaiknya
            berupa vektor kolom. Jika digunakan vektor baris atau matriks
            warna penuh, warna diterapkan pada tiap titik data.
thickness : ketebalan garis untuk kurva
            Nilainya dapat kurang dari 1 untuk garis yang sangat tipis.
style     : Gaya plot untuk garis, penanda titik, dan isian.
            Untuk titik gunakan:
            "[]", "<>", ".", "..", "...", "*", "+", "|", "-", "o"
            "[]#", "<>#", "o#" (bentuk terisi)
            "[]w", "<>w", "ow" (tidak transparan)
            Untuk garis gunakan:
            "-", "--", "-.", ".", ".-.", "-.-", "->"
            Untuk poligon terisi atau diagram batang gunakan:
            "#", "#O", "O", "/", "\", "\/", "+", "|", "-", "t"
points    : memplot titik tunggal alih-alih segmen garis
addpoints : jika true, memplot segmen garis sekaligus titik
add       : menambahkan plot ke plot yang sudah ada
user      : mengaktifkan interaksi pengguna untuk fungsi
delta     : ukuran langkah untuk interaksi pengguna
bar       : diagram batang (x adalah batas interval, y adalah nilai interval)
histogram : memplot frekuensi x dalam n subinterval
distribution=n : memplot distribusi x dengan n subinterval
even      : menggunakan nilai bilangan bulat untuk histogram otomatis
steps     : memplot fungsi sebagai fungsi tangga (steps=1,2)
adaptive  : menggunakan plot adaptif (n adalah jumlah langkah minimum)
level     : memplot garis ketinggian dari fungsi implisit dua variabel
outline   : menggambar garis batas dari rentang level
\end{verbatim}

Jika nilai \texttt{level} berupa matriks $2 \times n$, rentang level akan digambar dengan warna dan gaya isian yang ditentukan. Jika \texttt{outline} bernilai true, batasnya digambar menggunakan warna kontur. Dengan fitur ini, daerah $f(x,y)$ di antara dua batas dapat ditandai.

\begin{verbatim}
hue       : menambahkan gradasi warna pada plot level untuk menunjukkan nilai fungsi
contour   : Menggunakan plot level dengan garis ketinggian otomatis
nc        : jumlah garis level otomatis
title     : judul plot (bawaan "")
xl, yl    : label untuk sumbu x dan sumbu y
smaller   : jika >0, menyediakan ruang lebih luas di sebelah kiri untuk label
vertical  : Mengaktifkan atau menonaktifkan label vertikal. Mengubah variabel
            global verticallabels secara lokal untuk satu plot. Nilai 1 hanya
            mengatur teks vertikal, nilai 2 menggunakan label angka vertikal
            pada sumbu y.
filled    : mengisi bagian dalam kurva
fillcolor : warna isian untuk diagram batang dan kurva terisi
outline   : garis batas untuk poligon terisi
logplot   : mengatur plot logaritmik
            1 = logplot pada sumbu y,
            2 = logplot pada sumbu xy,
            3 = logplot pada sumbu x
own       : String yang menunjuk ke rutin plot buatan sendiri. Dengan >user, Anda
            mendapatkan interaksi pengguna yang sama seperti pada plot2d.
            Rentang akan diatur sebelum setiap pemanggilan fungsi Anda.
maps      : memetakan ekspresi (0 lebih cepat), fungsi selalu dipetakan.
contourcolor : warna garis kontur
contourwidth : ketebalan garis kontur
clipping  : mengaktifkan/menonaktifkan pemotongan (clipping) plot (bawaan true)
title     : Digunakan untuk mendeskripsikan plot dan muncul di bagian atas.
            Selain itu, label sumbu x dan y dapat ditambahkan dengan xl="string"
            atau yl="string". Label lain dapat ditambahkan dengan fungsi label()
            atau labelbox(). Judul dapat berupa string Unicode atau gambar
            rumus LaTeX.
cgrid     : Menentukan jumlah garis kisi untuk plot kisi kompleks.
            Sebaiknya merupakan pembagi dari (ukuran matriks dikurangi 1).
            cgrid dapat berupa vektor [cx,cy].
\end{verbatim}

\subsection*{Ringkasan}

Fungsi ini dapat memplot:
\begin{itemize}
    \item ekspresi, koleksi pemanggilan, atau fungsi satu variabel,
    \item kurva parametrik,
    \item data $x$ terhadap data $y$,
    \item fungsi implisit,
    \item diagram batang (bar plots),
    \item kisi bilangan kompleks,
    \item poligon.
\end{itemize}

Jika fungsi atau ekspresi diberikan untuk \texttt{xv}, \texttt{plot2d()} akan menghitung nilai-nilai pada rentang yang ditentukan menggunakan fungsi atau ekspresi tersebut. Ekspresi harus dinyatakan dalam variabel $x$. Rentang harus ditentukan melalui parameter \texttt{a} dan \texttt{b} kecuali jika menggunakan rentang bawaan $[-2,2]$. Rentang $y$ akan dihitung secara otomatis, kecuali jika \texttt{c} dan \texttt{d} ditentukan, atau jari-jari \texttt{r} diberikan sehingga menghasilkan rentang $[-r,r]$ untuk $x$ dan $y$. Untuk plot fungsi, \texttt{plot2d} menggunakan evaluasi adaptif secara bawaan. Untuk mempercepat pemrosesan fungsi yang rumit, matikan opsi ini dengan \texttt{<adaptive}, dan kurangi jumlah interval \texttt{n} bila perlu. Selain itu, \texttt{plot2d()} secara bawaan melakukan pemetaan elemen demi elemen. Jika ekspresi atau fungsi Anda dapat menangani vektor $x$ secara langsung, matikan pemetaan dengan \texttt{<maps} untuk evaluasi yang lebih cepat. Perhatikan bahwa plot adaptif selalu dihitung elemen demi elemen.

Jika fungsi atau ekspresi diberikan untuk \texttt{xv} sekaligus \texttt{yv}, \texttt{plot2d()} akan menggambar kurva dengan nilai \texttt{xv} sebagai koordinat $x$ dan nilai \texttt{yv} sebagai koordinat $y$. Dalam kasus ini, rentang parameter sebaiknya ditentukan menggunakan \texttt{xmin} dan \texttt{xmax}. Ekspresi yang ditulis di dalam string harus selalu dinyatakan dalam variabel parameter $x$.

\end{document}